\section{Introduction}\label{sec:intro}

\subsection{The result}
Let $\cX$ be a connected complex manifold of dimension four with a
K\"ahler form $\omega$, let $\DD\subset\CC$ be a disc centred at the
origin, and let
\[
 \pi:\cX\longrightarrow\DD
\]
be a proper surjective holomorphic map with connected fibres
$X_t=\pi^{-1}(t)$, which is a submersion over $\DD\setminus\{0\}$. We
assume that the fibre over the origin is reduced and that $X_0$ has
exactly $N\ge1$ singular points
$p_1,\ldots,p_N$, each an ordinary double point. For $t\ne0$ let $g_t$ be
the K\"ahler metric induced by $\omega$ on the smooth threefold $X_t$,
and let
\[
 0<\lambda^+_{2,1}(t)\le\lambda^+_{2,2}(t)\le\cdots
\]
be the positive eigenvalues of its Hodge Laplacian $dd^*+d^*d$ on real
two-forms, with multiplicity.

Let $q:Y\to X_0$ be a small resolution: $Y$ is a complex manifold, $q$ is
an isomorphism over $X_0\setminus\{p_1,\ldots,p_N\}$, and the
$C_i=q^{-1}(p_i)$ are smooth rational curves. Small resolutions exist as complex manifolds, since each node has two
local ones; $Y$ need not be K\"ahler, and the hypotheses below do not
require it. We assume
\begin{equation}\label{eq:hypotheses}
 b_1(Y)=0,\qquad H^2(Y;\RR)\ \text{is spanned by first Chern classes of
 holomorphic line bundles.}
\end{equation}
Both conditions hold, for instance, when $Y$ is projective with
$H^1(Y,\OO_Y)=H^2(Y,\OO_Y)=0$. Put $V=H^2(Y;\RR)$ and
$H=[q^*\omega]\in V$, and define
\begin{equation}\label{eq:intro-data}
\begin{split}
 \per&:V\longrightarrow\RR^N,\qquad \per(D)=(D\cdot C_i)_i,\qquad K=\ker\per,\\
 G(D,E)&=-\int_YD\,E\,H+\frac{3\,(\int_YD\,H^2)(\int_YE\,H^2)}
 {2\int_YH^3}.
\end{split}
\end{equation}
The vector $\per(D)$ records the degrees of $D$ on the exceptional
curves. The kernel $K$ is the image of $H^2(X_t;\RR)$ under restriction to the
complement of small neighbourhoods of the nodes, which is common to all
fibres, and it contains $H$ (Lemma~\ref{lem:topology}). We show that $G$ is
positive definite; it is the $L^2$ inner product of the harmonic
representatives on the regular part of $X_0$, and on
$\{\int_YDH^2=0\}$ it agrees with the Hodge--Riemann form
$-\int_YD\,E\,H$ of $H$. These classes are $G$-orthogonal to $H$, and
$G(H,H)=\tfrac12\int_YH^3$. Let $\cM$
be the form induced by $G$ on $V/K$, that is, the restriction of $G$ to
the $G$-orthogonal complement of $K$; equivalently
\[
 \cM([D],[D])=\min_{k\in K}G(D+k,D+k).
\]
Let $\cE([D],[E])=\per(D)\cdot\per(E)$, and let
\[
 0<\mu_1\le\cdots\le\mu_m,\qquad m=\dim V/K=b_2(Y)-b_2(X_t),
\]
be the eigenvalues of $\cE$ relative to $\cM$, that is, the roots of
$\det(\cE-\mu\cM)=0$. They do not depend on the choice of small
resolution (Remark~\ref{rem:G}).

\begin{theorem}\label{thm:main}
Under the preceding hypotheses,
\begin{equation}\label{eq:main}
 \lim_{t\to0}\,\log(1/|t|)\,\lambda^+_{2,j}(t)=2\pi\mu_j
 \qquad(1\le j\le m).
\end{equation}
There is $\gamma>0$ such that, for all sufficiently small $t\ne0$, the
positive eigenvalues of the Hodge Laplacian on one-forms and the
eigenvalues $\lambda^+_{2,j}(t)$ with $j>m$ are at least $\gamma$. For
such $t$ the eigenforms for $\lambda^+_{2,1}(t),\ldots,\lambda^+_{2,m}(t)$
are coexact, primitive and of type $(1,1)$.
\end{theorem}

Thus the number of small eigenvalues is the number of independent classes
in $H^2(Y)$ that do not extend to the smooth fibres, and the limits are
computed from the intersection form of $Y$, the class of the metric and
the degrees of these classes on the exceptional curves. When $m=1$ and
$K=\RR H$,
\begin{equation}\label{eq:intro-rank-one}
 \lambda^+_{2,1}(t)=\frac{2\pi\sum_i(D\cdot C_i)^2}
 {(\int_YDH^2)^2/\int_YH^3-\int_YD^2H}\ \frac{1+o(1)}{\log(1/|t|)}
\end{equation}
for any class $D$ independent of $H$. For the determinantal hypersurface
$\{l_1f_2-l_2f_1=0\}\subset\PP^4$ of degree $d\ge2$, with $(d-1)^2$
nodes, smoothed in a pencil and given the Fubini--Study metric with lines
of area one, the coefficient is $2\pi d$
(Example~\ref{ex:determinantal}).

\subsection{Relation to earlier work}
Ann\'e and Takahashi \cite{anne-takahashi} study the Hodge--de Rham
operator on a manifold in which a compact piece with conical boundary is
scaled to zero. They determine the multiplicity of zero in the limit
spectrum and obtain convergence of the remaining spectrum; in degrees
related to the cohomology of the link they construct test forms of
energy $O(1/|\log\epsilon|)$ in the scaling parameter $\epsilon$, which
bounds certain eigenvalues above. We use their compactness results
\cite[Corollary~17 and Lemma~18]{anne-takahashi}. The leading constants of such eigenvalues are not determined there.

In \cite{wuebben-conifold} the limits \eqref{eq:main} are proved, along
rays $t\to0$, for the Ricci-flat K\"ahler metrics in a polarized
Calabi--Yau conifold transition, with the same forms $\cE$ and $\cM$.
That proof uses the description of the Ricci-flat metrics near the
singular fibre due to Hein and Sun \cite{hein-sun} and Benabida
\cite{benabida}. For the induced metrics considered here the geometry
near the nodes is explicit (Lemma~\ref{lem:normal-form}), no condition
on the canonical bundle is needed, and the limit is taken over all
$t\to0$, not only along rays. The tangent cone at a node, its link, the local smoothing of the cone,
and the constants attached to them differ from the Ricci-flat case, while the
limits agree whenever both theorems apply
(Remark~\ref{rem:metric}); the limits depend on the K\"ahler metric only
through its class, because the squared norm of the harmonic form on the
link and the logarithm of the scale of the local smoothing change by the
same factor (Remark~\ref{rem:metric}). The statement about types is proved differently
here, without a parallel holomorphic volume form
(Proposition~\ref{prop:type}).

The spectral argument of Section~\ref{sec:spectrum} adapts that of
\cite[Sections~4--6]{wuebben-conifold} to the cone $\{\sum_j\sigma_jw_j^2=0\}\subset\CC^4$ of
Section~\ref{sec:node}; we
give it in full, so that no result of \cite{wuebben-conifold} is used.
New here are the models of Section~\ref{sec:node},
the construction of the harmonic representatives without a Ricci-flat
background (Proposition~\ref{prop:harmonic}), the passage to arbitrary
$t\to0$ (Lemma~\ref{lem:condition-N}), and the proof of the statement
about types.

For functions, small eigenvalues of metrics induced on degenerating
families are studied by Dai and Yoshikawa \cite{dai-yoshikawa} for curves
and by Cao \cite{cao} in all dimensions; they have order
$1/\log(1/|t|)$, and for a reduced central fibre their number is the
number of its components minus one
\cite[Theorem~0.2]{dai-yoshikawa}, \cite[Theorem~0.4]{cao}. Dai and Yoshikawa raise the question
of small eigenvalues for forms of higher degree: for $(n,0)$-forms with
induced metrics and for forms of type $(p,q)$ with K\"ahler--Einstein
metrics of negative scalar curvature
\cite[Conjecture~9.5 and Problem~9.7]{dai-yoshikawa}.
Theorem~\ref{thm:main} is a result of this kind in a case not covered by
those formulations, two-forms of type $(1,1)$ with induced metrics
on threefolds acquiring ordinary double points; here the central fibre
is irreducible and there are no small eigenvalues on functions.

\subsection{Method}
In suitable holomorphic coordinates at a node the map $\pi$ is a
diagonal quadratic form $\sum_j\sigma_jw_j^2$ and $\omega$ is the flat
K\"ahler form up to terms vanishing at the node. For the flat form, the
fibres of $\pi$ near the node are exact rescalings, by $|t|^{1/2}$, of
the affine quadric $\{\sum_j\sigma_jw_j^2=1\}$, and the central fibre is
the cone over a circle bundle $L_\sigma$ over a quadric surface. The
harmonic two-form on $L_\sigma$ representing a generator of
$H^2(L_\sigma;\ZZ)$ has squared $L^2$ norm $\pi$, independently of $\sigma$ (Lemma~\ref{lem:link-form}).

A class $D\in V$ has a closed and coclosed representative on the regular
part of $X_0$, asymptotic to $(D\cdot C_i)$ times that form at the
$i$th node (Proposition~\ref{prop:harmonic}). It extends to a closed
form on $X_t$ only if all these periods vanish. Otherwise, cutting it
off between the scale $|t|^{1/2}$ of the rescaled affine quadric and a
fixed scale costs energy, that is $\|d\psi\|^2+\|d^*\psi\|^2$, equal to
$\pi\sum_i(D\cdot C_i)^2/\log(1/|t|^{1/2})$ to leading order, and a lower bound for the energy on the affine quadric shows that no form with the same periods
has less (Lemma~\ref{lem:cap-energy}). The limits \eqref{eq:main} follow
from this on the model manifolds of Definition~\ref{def:model}, which
are exactly conical near the nodes (Theorem~\ref{thm:model-law}), using the compactness theorem of
\cite{anne-takahashi}, and are transferred to the metrics $g_t$ by
Dodziuk's comparison of Hodge eigenvalues for uniformly equivalent
metrics \cite{dodziuk}, which requires only that $g_t$ be close to the
models in $C^0$ (Lemma~\ref{lem:condition-N}). The statement about types
follows because the projection onto forms of type $(2,0)+(0,2)$ commutes
with the Laplacian of a K\"ahler metric, while the limits of the small
eigenforms are of type $(1,1)$.

\subsection{Organization}
Section~\ref{sec:node} contains the local geometry and the models,
Section~\ref{sec:harmonic} the harmonic representatives and the form $G$,
Section~\ref{sec:spectrum} the spectral analysis of the models, and
Section~\ref{sec:proof} the proof of Theorem~\ref{thm:main}.
Section~\ref{sec:examples} contains the example and remarks.
